% Generated draft: requires mathematical and source review.
\documentclass[11pt]{article}
\usepackage[margin=25mm]{geometry}
\usepackage{fontspec}
\setmainfont{Latin Modern Roman}
\usepackage{amsmath,amssymb,mathtools,unicode-math}
\setmathfont{Latin Modern Math}
\usepackage{enumitem,needspace,xurl,hyperref}
\setlength{\parindent}{0pt}
\setlength{\parskip}{5pt}
\setlength{\emergencystretch}{3em}
\newcommand{\N}{\mathbb{N}}
\newcommand{\Z}{\mathbb{Z}}
\newcommand{\Q}{\mathbb{Q}}
\newcommand{\R}{\mathbb{R}}
\newcommand{\C}{\mathbb{C}}
\newcommand{\F}{\mathbb{F}}
\newcommand{\A}{\mathbb{A}}
\newcommand{\PP}{\mathbb{P}}
\newcommand{\E}{\mathbb{E}}
\newcommand{\eps}{\varepsilon}
\newcommand{\dd}{\,\mathrm{d}}
\begin{document}
\section*{UnsolvedMath ID: 20002360}
\textbf{Source problem code:} catalogue entry
\subsection{Definitions and mathematical statement}
Let \(K=\C(Z)\) be the rational-function field in one indeterminate \(Z\), not the complex field expanded by a predicate for the integers. For \(a,b\in\C\) with \(4a^3+27b^2\ne0\), let
\[
 E_{a,b}:Y^2=X^3+aX+b,\qquad
 j(E_{a,b})=1728\frac{4a^3}{4a^3+27b^2}.
\]
Write \(\operatorname{End}_{\C}(E)\) for the endomorphism ring of the elliptic curve with its identity. In characteristic zero its underlying abelian group has rank one or two. Complex multiplication means \(\operatorname{End}_{\C}(E)\ne\Z\), where \(\Z\) denotes the scalar multiplication endomorphisms. Put
\[
 S_{\mathrm{CM}}=\{j(E)\in\C:E/\C\text{ has complex multiplication}\}.
\]
These are the singular moduli. The source uses the criterion
\[
 |\operatorname{End}(E)/2\operatorname{End}(E)|>2
 \quad\Longleftrightarrow\quad \operatorname{End}(E)\ne\Z.
\]
The rank-one/rank-two fact for complex elliptic curves is essential; freeness and finite generation alone would not imply this for an arbitrary ring or module.

\textbf{Question 23 (Pheidas).} The source reports that the displayed criterion is Diophantine and that \(S_{\mathrm{CM}}\), viewed as a subset of \(K\), can be existentially defined in \(K\). What arithmetic or logical consequences can be obtained from this definable set? The question does not specify a desired consequence. In particular, a Diophantine interpretation of integer arithmetic would require additional definable data and operations, not just a unary predicate. The claimed existential formula and its permitted parameters are not supplied, so existential definability is retained as a source assertion rather than an explicit construction in this draft.

\textbf{Reported certificate obstructions.} If the endomorphism order of a CM elliptic curve has discriminant \(D<0\), the aid reports that every non-scalar \(\alpha\in\operatorname{End}(E)\setminus\Z\) satisfies
\[
 \deg(\alpha)\geq\frac{|D|}{4}.
\]
Here degree is the degree of the elliptic-curve morphism, distinct from rational-function witness degree below. Consequently, for every fixed bound \(N\), singular moduli admitting a non-scalar endomorphism of degree at most \(N\) form a finite set. No converse based only on this inequality is asserted.

More generally, fix a Diophantine formula \(\Phi(t)\) over \(K\), with finitely many fixed coefficients and a fixed finite tuple of existential witnesses, such that
\[
 \{c\in\C:K\models\Phi(c)\}=S_{\mathrm{CM}}.
\]
For a rational function \(h=A(Z)/B(Z)\) in reduced form, \(B\ne0\), use \(\deg_Z(h)=\max\{\deg A,\deg B\}\), with constants, including zero, of degree zero. For each integer \(d\geq0\), let
\[
 S_{\Phi,d}=\{c\in S_{\mathrm{CM}}:\Phi(c)\text{ has witnesses in }K
 \text{ each of }Z\text{-degree at most }d\}.
\]
The supplied aid reports
\[
 |S_{\Phi,d}|<\infty\quad\text{for every fixed }\Phi,d,
 \qquad S_{\mathrm{CM}}=\bigcup_{d\geq0}S_{\Phi,d}.
\]
Thus, for a fixed formula, no uniform finite witness-degree bound covers the entire CM locus. The formula, its coefficients and its witness presentation must remain fixed as \(c\) varies. This statement is conditional on an exact Diophantine definition on constants and is not a non-definability theorem. It does not rule out unbounded certificates, additional definable relations, or an eventual interpretation of arithmetic. The record supplies no such interpretation or decision-theoretic consequence.
\subsection{Short English statement}
Pheidas asks what can be deduced from the reported existential definition of CM \(j\)-invariants in \(\C(Z)\). The aid states that bounded-degree endomorphisms, and bounded rational-function witnesses for any fixed definition, account for only finitely many singular moduli. This requires unbounded certificates and leaves the requested arithmetic applications unspecified.
\subsection{Sources}
\begin{itemize}
\item Local UnsolvedMath registry, permanent ID \(20002360\), catalogue entry, \texttt{aim-logic-notes.json}, \texttt{notes[135]}, workshop \emph{Extensions of Hilbert's tenth problem}, Question 23 (Pheidas): \url{https://aimath.org/WWN/hilberts10th/hilberts10th.pdf}.
\item Supplied research summary on CM discriminants, endomorphism-degree bounds and bounded rational-function witness strata. No separate authoritative citation or explicit Diophantine formula is supplied.
\item Attempt 1 is labelled partial\_result, problem status at run ``open'', rigor ``proved'', low novelty confidence; the registry currently says partially\_solved. These labels concern a machine-generated research aid, not peer-reviewed verification; independently verify before citation. The supplied statement-verification field reports inspection of page 5 of the workshop PDF; this draft has not independently visited it.
\end{itemize}
\end{document}
